\HMTNarrativeHeading{De recuperar un estado a transportar una familia}

La recuperación compatible permite extender la exposición desde
un estado hasta familias de estados. Una lectura enriquecida
individualiza sus antecedentes cuando su memoria distingue los
elementos que comparten el mismo valor visible. La inversa
restituye entonces el estado sobre la imagen de esa lectura.
La compatibilidad con el refinamiento reúne las inversas finitas
en una reconstrucción de historias completas.

El transporte de familias exige conservar la pertenencia. Una
biyección de historias lleva cada familia a su imagen y cada
elemento a su correspondiente. Las uniones, intersecciones y
complementos relativos se transportan mediante esas operaciones.
La restricción de la biyección a una familia conserva su cardinal.
La teoría de memoria adquiere así un enlace preciso con las
operaciones sobre conjuntos de historias.

La medida requiere además el transporte de los pesos. Una
simetría compatible con el árbol de refinamiento conserva la
medida de cilindros cuando respeta las emisiones y sus
multiplicidades. La naturalidad expresa la concordancia entre
transportar y truncar. Las demostraciones siguientes explicitan
estas relaciones sobre los objetos definidos en el artículo.

El estudio específico de la hipótesis del continuo desarrolla
la clausura genealógica y su jerarquía de conjuntos. Su resultado
cardinal se formula sobre esa clausura y su potencia interna.
La relación que esta parte incorpora es la recuperación de
estados y el transporte de familias que mantienen la estructura
del continuo generado. El enunciado cardinal conserva su
desarrollo especializado y su dominio expresamente identificado.

Esta articulación sitúa el continuo dentro del argumento de
conservación. La historia completa, la familia a la que pertenece
y la medida de su cilindro requieren operaciones relacionadas,
con contenidos diferentes. El artículo reúne sus condiciones de
compatibilidad para mostrar cómo la doble proyección se mantiene
al pasar de una lectura individual a la organización de familias.
La unidad de la exposición reside en esos mapas y en las pruebas
que preservan sus relaciones.

